\section{Marco Teórico}

El presente marco teórico reúne los conceptos necesarios para comprender el diseño de
una plataforma restaurantera orientada a datos. Aunque el sistema incluye interfaces para
comensales, meseros, cocina y administración, su propósito técnico central es capturar
eventos operativos de manera estructurada, almacenarlos en una base de datos transaccional,
garantizar su trazabilidad y transformarlos posteriormente en indicadores analíticos mediante
un pipeline batch con Apache Spark.

Por ello, los fundamentos teóricos se organizan en cinco áreas: el dominio operativo
restaurantero, los sistemas transaccionales y bases de datos relacionales, la gobernanza y
seguridad de los datos, la sincronización en tiempo real entre actores, y la arquitectura
analítica basada en procesamiento por lotes, modelo Medallion e inteligencia de negocio.

\subsection{Dominio de Aplicación: Operación Restaurantera}

\subsubsection{Restaurantes como Entornos de Operación y Servicio}

Los restaurantes constituyen entornos de operaci\'on y servicio en los que convergen
procesos de atenci\'on al cliente, coordinaci\'on del personal, gesti\'on de \'ordenes,
preparaci\'on de alimentos y registro de consumos. Dentro de este entorno suelen
distinguirse dos grandes \'ambitos operativos: el \textit{front of house} (contacto directo
con el cliente) y el \textit{back of house} (tareas internas de preparaci\'on y
coordinaci\'on) \cite{lee_digital_2024}. La digitalizaci\'on de estos procesos busca
mejorar la velocidad, precisi\'on y conveniencia operativa \cite{nra_state_2023}.

\subsubsection{Front of House, Back of House y KDS}

La coordinaci\'on entre el sal\'on y la cocina es cr\'itica para la eficiencia del servicio.
Un Sistema de Pantallas de Cocina (KDS, \textit{Kitchen Display System}) es una interfaz
digital utilizada para visualizar, organizar y actualizar órdenes dentro del área de cocina.
Sustituye o complementa las comandas físicas, permitiendo que los pedidos se agrupen por
estado, prioridad o estación de preparación.

En el presente proyecto, el KDS no solo cumple una función operativa, sino también una
función de captura de datos. Cada transición de estado de una orden, como pendiente, en
preparación, lista o entregada, puede registrarse con marcas de tiempo. Estos eventos
permiten calcular métricas como duración de preparación, carga de trabajo por estación y
tiempos promedio de atención, las cuales posteriormente alimentan el pipeline analítico.

\subsubsection{Omnicanalidad en Sistemas Digitales de Servicio}

La omnicanalidad se entiende como un enfoque integrado en el que múltiples canales o
interfaces de interacción operan sobre un flujo de información común, evitando que cada
punto de contacto funcione como un sistema aislado \cite{wang_omnichannel_2024}. En
contextos de servicio, esto implica coordinar datos, procesos y estados entre actores
distintos.

Para efectos del presente trabajo, la omnicanalidad se define como la integración
coordinada de las interfaces del comensal, mesero, cocina y administración dentro de un
mismo flujo operativo. Cada actor interactúa con una vista diferente del sistema, pero las
acciones registradas alimentan una misma base de datos transaccional. De esta forma, una
orden creada desde la mesa virtual puede reflejarse en cocina, actualizar el estado de la
mesa y posteriormente formar parte del pipeline analítico.

\subsubsection{Aplicaciones Web Progresivas (PWA)}

Las Aplicaciones Web Progresivas (PWA) son aplicaciones construidas con tecnologías web
que buscan ofrecer una experiencia similar a la de una aplicación instalada. De acuerdo
con MDN Web Docs, una PWA puede ejecutarse desde una sola base de código en distintos
dispositivos y ofrecer capacidades como instalación, ejecución en navegador, uso de
\textit{web app manifest} e integración parcial con el sistema operativo \cite{mdn_pwa, mdn_manifest}.

Una PWA no implica necesariamente que todas sus funcionalidades estén disponibles sin
conexión. Aunque los \textit{service workers} permiten implementar estrategias de caché y
experiencias offline, en sistemas transaccionales colaborativos puede ser necesario exigir
conectividad permanente para mantener la consistencia de la información.

En el presente proyecto, la PWA se utiliza como mecanismo de acceso ligero para comensales
y personal operativo. Su propósito es permitir el uso desde dispositivos móviles sin
instalación obligatoria desde tiendas de aplicaciones, facilitar el acceso mediante código QR
y mantener una experiencia similar a una aplicación instalada. Debido a que las órdenes,
estados de cocina y liquidaciones internas requieren sincronización en tiempo real, la
versión 1.0 no contempla operación offline.

\subsubsection{Códigos QR como Mecanismo de Acceso a Mesas Virtuales}

Los códigos QR son códigos bidimensionales capaces de almacenar información de manera
compacta y ser leídos rápidamente por dispositivos móviles. Su estandarización bajo
ISO/IEC 18004 ha favorecido su uso como mecanismo de acceso a recursos digitales en
contextos comerciales y de servicio \cite{denso_qr_standards}.

En restaurantes, los códigos QR se utilizan comúnmente para consultar menús digitales,
iniciar pedidos desde la mesa o acceder a servicios asociados al establecimiento. En el
presente proyecto, el QR funciona como punto de entrada a una mesa virtual. El enlace
contiene un identificador de mesa acompañado de una firma HMAC-SHA256 generada en el
servidor, lo que permite validar que el identificador no haya sido manipulado.

Una vez validado el enlace, el acceso a la mesa depende del estado de la sesión activa y
de la clave definida para la mesa virtual. De esta manera, el QR no se plantea como un
mecanismo de autenticación completo, sino como un mecanismo de acceso inicial asociado a
una entidad física del restaurante.

\subsection{Fundamentos de Sistemas Transaccionales}

\subsubsection{Bases de Datos Relacionales y Propiedades ACID}

Las bases de datos relacionales permiten representar información estructurada mediante
tablas, atributos, claves primarias, claves foráneas y relaciones entre entidades. Este
modelo es especialmente adecuado para sistemas transaccionales donde se requiere mantener
integridad, consistencia y trazabilidad entre operaciones relacionadas \cite{ibm_relational_db}.

Las propiedades ACID describen garantías fundamentales para este tipo de sistemas:
atomicidad, consistencia, aislamiento y durabilidad. La atomicidad asegura que una
transacción se ejecute completamente o no se ejecute; la consistencia mantiene las reglas
definidas por el modelo de datos; el aislamiento evita interferencias indebidas entre
transacciones concurrentes; y la durabilidad garantiza la persistencia de los cambios
confirmados.

En el presente proyecto, la base de datos relacional representa el núcleo transaccional
del sistema. Entidades como cadena, zona, sucursal, usuarios, mesas, sesiones, comensales,
productos, órdenes, estados de cocina y liquidaciones internas se modelan de forma
relacionada para permitir trazabilidad completa desde el evento operativo hasta su uso
analítico posterior.

\subsubsection{Integridad Referencial y Snapshots Históricos}

En sistemas transaccionales, no basta con almacenar referencias a entidades vivas del
catálogo, ya que estas pueden modificarse con el tiempo. Por ejemplo, el nombre, precio,
categoría o modificadores de un producto pueden cambiar después de que una orden haya sido
registrada. Si una orden histórica dependiera únicamente del estado actual del catálogo,
su interpretación podría alterarse.

Para evitar este problema se utilizan snapshots históricos, es decir, copias de los datos
relevantes al momento de la transacción. En el contexto del presente proyecto, cada ítem de
orden debe conservar información como nombre del producto, precio aplicado, variante,
modificadores y cantidad al momento de la compra. Esto permite que las órdenes históricas
mantengan su significado original, incluso si el catálogo cambia posteriormente.

Este principio es fundamental para la calidad de los datos analíticos, ya que las capas
Silver y Gold del pipeline dependen de registros históricos consistentes. Sin snapshots,
métricas como ingresos por producto, consumo por categoría o comportamiento de demanda
podrían distorsionarse al reflejar valores actuales en lugar de valores vigentes al momento
de la operación.

\subsubsection{Liquidación Interna y Consistencia Transaccional}

La liquidación interna se refiere al registro lógico del estado de pago o aportación
dentro del sistema, sin que necesariamente exista procesamiento bancario real. En el caso
de una mesa de restaurante, distintos comensales pueden cubrir consumos individuales,
aportar a una cuenta compartida o completar el total pendiente de la sesión. Estas acciones
requieren validaciones para evitar inconsistencias en los montos registrados.

En la versión 1.0 del presente proyecto, la liquidación no procesa tarjetas, transferencias
ni billeteras digitales. El sistema únicamente registra aportaciones internas capturadas
por los usuarios y actualiza el estado de la cuenta con base en los consumos registrados.
Por ello, el problema principal no es la integración financiera, sino la consistencia de los
datos transaccionales.

Cuando dos o más comensales intentan registrar aportaciones simultáneamente sobre una misma
mesa, pueden presentarse condiciones de carrera. Para prevenir que el monto liquidado exceda
el total pendiente, se pueden utilizar mecanismos de bloqueo a nivel de base de datos, como
\texttt{SELECT FOR UPDATE} en PostgreSQL. Este tipo de control permite leer y bloquear la fila
correspondiente durante una transacción, validar el estado actual y confirmar únicamente las
operaciones consistentes.

\subsubsection{Comunicación en Tiempo Real mediante WebSockets y Supabase Realtime}

La comunicación en tiempo real permite que múltiples clientes conectados reciban cambios
de estado sin necesidad de recargar la página o consultar continuamente al servidor. La API
de WebSockets establece un canal bidireccional persistente entre cliente y servidor, útil
para escenarios donde los eventos deben propagarse con baja latencia \cite{mdn_websocket}.

En el contexto restaurantero, esta capacidad permite sincronizar acciones entre comensales,
meseros, cocina y administración. Por ejemplo, una orden registrada desde la mesa virtual
puede aparecer en el tablero de cocina, y un cambio de estado realizado por el personal de
cocina puede reflejarse en la vista del comensal.

Supabase Realtime utiliza conexiones WebSocket para entregar actualizaciones a clientes
suscritos a eventos o cambios específicos. En el presente proyecto, esta funcionalidad se
emplea como mecanismo de sincronización operativa para órdenes, estados de mesa y estados
de preparación.

\subsection{Seguridad y Gobernanza de Datos}

\subsubsection{Gobernanza de Datos}

La gobernanza de datos se define como la ejecución de autoridad y control sobre la gestión
de los activos de datos. Establece el marco estratégico para asegurar la disponibilidad,
usabilidad, integridad y seguridad de la información organizacional \cite{dama2017dmbok}.
En este proyecto, la gobernanza asegura que los datos transaccionales recopilados mantengan
un nivel de calidad óptimo para su posterior explotación analítica.

\subsubsection{RBAC y RLS}

El Control de Acceso Basado en Roles (RBAC) restringe el acceso basándose en las funciones
de los usuarios \cite{sandhu1996role}. En el presente proyecto, RBAC define qué módulos
puede utilizar cada rol: Admin Cadena, Gerente Zona, Gerente Sucursal, Mesero, Cocina/Chef
y Comensal.

La Seguridad a Nivel de Fila (RLS) complementa este modelo desde la base de datos, asegurando
que cada usuario solo consulte o modifique registros correspondientes a su cadena, zona,
sucursal o sesión de mesa \cite{postgresql_rls}. Esta combinación reduce el riesgo de
depender únicamente de validaciones en la capa de aplicación.

\subsubsection{Seguridad en Tránsito mediante HTTPS}

La seguridad de las comunicaciones se garantiza mediante el protocolo HTTPS y TLS 1.2+, que
proporcionan confidencialidad, integridad y autenticación mediante el cifrado de los datos
en tránsito \cite{rescorla2018tls}. Esto protege la transmisión de credenciales de acceso,
identificadores de mesa y telemetría operativa, asegurando que la información no sea
interceptada ni alterada.

\subsection{Arquitectura Analítica}

\subsubsection{OLTP y OLAP}

En arquitectura de datos se distingue entre sistemas OLTP (\textit{Online Transaction
Processing}) y sistemas OLAP (\textit{Online Analytical Processing}). Los sistemas OLTP están
optimizados para registrar transacciones frecuentes, mantener consistencia y responder
rápidamente a operaciones individuales. En este proyecto, la base de datos transaccional
registra eventos como órdenes, cambios de estado y liquidaciones internas.

Los sistemas OLAP, en cambio, están orientados al análisis histórico, consultas agregadas
y exploración de grandes volúmenes de datos. En lugar de consultar directamente la base
operativa para generar reportes complejos, es común construir una capa analítica separada
que reduzca la carga sobre el sistema transaccional y permita organizar los datos para
inteligencia de negocio.

Esta separación resulta relevante para el presente proyecto, ya que permite mantener la
operación restaurantera en tiempo real sobre una base OLTP, mientras que los indicadores
administrativos se generan a partir de datasets analíticos procesados en lote.

\subsubsection{Procesamiento por Lotes y Apache Spark}

El procesamiento por lotes (\textit{batch processing}) consiste en ejecutar tareas sobre
conjuntos de datos acumulados previamente, sin requerir respuesta inmediata para el usuario
final. Este enfoque es adecuado para generar reportes, limpiar datos, calcular agregados y
preparar información para análisis histórico \cite{kleppmann2017designing}.

Apache Spark es un motor de procesamiento utilizado en ingeniería de datos para ejecutar
transformaciones sobre grandes conjuntos de información. Su modelo permite trabajar con
datos distribuidos, realizar operaciones de limpieza, agregación y enriquecimiento, y
construir pipelines analíticos repetibles \cite{zaharia2012resilient, chambers2018spark}.

En el presente proyecto, Spark se emplea para procesar por lotes los datos transaccionales
generados por la operación restaurantera. Aunque la validación inicial se realizará en un
entorno controlado (VPS), el uso de Spark permite diseñar una arquitectura preparada para
escalar cuando aumente el volumen de datos. Los jobs desarrollados transformarán registros
de mesas, órdenes, productos, estados de cocina y liquidaciones internas en capas analíticas
Bronze, Silver y Gold.

\subsubsection{Arquitectura Medallion: Bronze, Silver y Gold}

La arquitectura Medallion es un patrón de ingeniería de datos que organiza el flujo de
información en capas progresivas de refinamiento. Su objetivo es mejorar la calidad,
estructura y utilidad de los datos conforme avanzan desde su estado original hasta su forma
analítica final \cite{armbrust2021lakehouse, reis2022fundamentals}.

La capa Bronze almacena los datos crudos extraídos desde los sistemas origen. En el presente
proyecto, esta capa conserva registros provenientes de la base transaccional, como órdenes,
productos, mesas, comensales, estados de cocina y liquidaciones internas.

La capa Silver contiene datos limpios, normalizados y enriquecidos. En esta etapa se corrigen
formatos, se resuelven relaciones entre entidades, se calculan duraciones operativas y se
preparan los datos para análisis posterior.

La capa Gold representa la capa de consumo analítico. En ella se generan agregados orientados
a negocio, como ventas por sucursal, productos más consumidos, tiempos promedio de preparación,
actividad por mesa y comportamiento por franja horaria.

\subsubsection{Inteligencia de Negocio y Dashboards}

La inteligencia de negocio (BI) permite convertir los datos procesados en información
estratégica para la toma de decisiones. Mediante dashboards administrativos, los gerentes
y administradores pueden visualizar el desempeño de la cadena restaurantera en sus distintos
niveles. En este proyecto, el dashboard consume directamente la capa Gold generada por el
pipeline de Spark, asegurando que las métricas reflejen la realidad operativa con alta
fidelidad y trazabilidad.
